\chapter{Results}
\label{ch:results}

\secref{results-foveation} reports the gaze-contingent input-foveation experiment specified in
\secref{foveation} and \secref{training}, the main experiment of this thesis, the seven arms of
\tabref{arms} over ten cross-validation folds.
\secref{results-per-image} reports the per-image diagnostic of \subsecref{difficulty-proxies}, which uses the
pretrained \dg{} read-out and characterises where the model and the human population agree.
The two sections use different read-outs on different splits, and \tabref{ledger} says which.

\begin{table}[tbp]
  \centering
  \caption[Each measurement, its read-out and its split]{Every measurement in this chapter, with the read-out it uses
    and the split it reads. The rows marked \emph{pretrained} take \dg{} as released and say nothing
    about the trained arms. Across the ten models each image is the test image exactly once and the
    validation image exactly once, so both splits span the dataset (\subsecref{cv-split}).}
  \label{tab:ledger}
  \footnotesize
  \setlength{\tabcolsep}{4pt}
  \begin{tabular}{lccc}
    \toprule
    measurement & read-out & split & sample \\
    \midrule
    Baseline (\subsecref{results-baseline})            & pretrained            & whole dataset & $104{,}171$ fixations \\
    Rate and reporting epoch (\subsecref{training-loop}) & each arm's own      & validation    & per fold \\
    Unadapted penalty (\subsecref{results-arms})       & pretrained            & validation    & $4$ arms \\
    Arm comparison (\subsecref{results-arms})          & each arm's, epoch $5$ & test          & $7$ arms \\
    Diagnostics (\subsecref{results-diagnostics})      & each arm's, epoch $5$ & validation    & $3$ contrasts \\
    Per-image variation (\secref{results-per-image})   & pretrained            & whole dataset & 1003 images \\
    \bottomrule
  \end{tabular}
\end{table}

\section{Gaze-contingent input foveation}
\label{sec:results-foveation}

\subsection{Baseline}
\label{subsec:results-baseline}

Pretrained \dg{} on the full MIT1003 dataset under the protocol of \subsecref{mit1003}
gains $1.562$~bits per fixation over the centerbias, pooled over fixations, which is
$0.026$~bits above the $1.536 \pm 0.010$ of \citet{kummerer2022deepgaze3}, Table~2. Its AUC (\eqnref{auc}) is
$0.9175 \pm 0.0018$ against their $0.916 \pm 0.001$, and its NSS
$3.265 \pm 0.066$ against their $3.257 \pm 0.016$. Both are per-image means over the 1003
stimuli.\footnote{Every $\pm$ in this chapter, and every error bar, is two standard errors of the mean
(\subsecref{stats-plan}), over the ten folds for every arm difference and over the 1003 images for the
pretrained baseline above. The values quoted from
\citet{kummerer2022deepgaze3} carry their convention rather than this one, a standard
deviation over eight repeated training runs.}
Each bit doubles the probability the model puts on the pixel a subject fixated (\subsecref{metrics}), so
$1.562$~bits is about $3$ times the probability the centerbias puts there.
This baseline is the state every arm's read-out starts from. Part of its margin over the published
figure is in-sample, since the released read-outs were fitted on these same images (\subsecref{arms});
the comparator for every difference below is the trained sharp control.

\subsection{The seven arms}
\label{subsec:results-arms}

Every arm fine-tunes its read-out on its own input, for the reason \secref{training} gives, and
everything below comes from those trained read-outs on the held-out test split.

Before that fine-tuning, the pretrained read-out on foveated input is level with the same read-out on
sharp input at the human cutoff, scoring $-0.0003 \pm 0.0034$~bits against it on validation, and falls
behind by $0.0224 \pm 0.0071$~bits at $f_{c0} = 20$ and by $0.0741 \pm 0.0123$ at $10$.
Read-out training recovers about three quarters of the first penalty and half of the second, measured
against the trained arms on the same split; at $f_{c0} = 40$ there was none to recover
(\figref{training-curves}).

\tabref{phase17-arms} reports the production run: all seven arms of \tabref{arms}, trained under the
protocol of \subsecref{training-loop} and read at epoch five on the held-out test folds, which together
cover the dataset once.
Every arm sits about $1.5$~bits above the centerbias, so no read-out collapsed.
The sharp control scores $1.5742$~bits, AUC $0.9180$ and NSS $3.309$, against the published $1.536 \pm 0.010$, $0.916 \pm 0.001$ and
$3.257 \pm 0.016$ of \subsecref{results-baseline}.
The fold-paired differences below agree with their validation counterparts within $0.0007$~bits on
every arm, so choosing the reporting epoch on
validation did not flatter that split.

% Generated by scripts/make_results_figures.py — do not edit by hand.
% Source: results/foveation_mit1003_initial/{test/table.json,
% test/per_image_ig.json}. Regenerate after any re-evaluation.
% Requires \usepackage{booktabs}.

\begin{table}[tbp]
\centering
\caption[Seven arms, ten folds]{Seven arms, ten image-stratified folds, test split at epoch 5.
The first three columns are means over the ten folds; the last pairs the folds and carries two standard
errors of its mean.}
\label{tab:phase17-arms}
\small
\setlength{\tabcolsep}{4.5pt}
\begin{tabular}{lcccr}
\toprule
arm & IG over centerbias & NSS & AUC & $\Delta$IG vs control, fold-paired \\
\midrule
sharp control & $1.5742$ & $3.309$ & $0.9180$ & --- \\
gaze-contingent @ 40 & $1.5781$ & $3.313$ & $0.9183$ & $+0.0039 \pm 0.0032$ \\
gaze-contingent @ 20 & $1.5684$ & $3.302$ & $0.9177$ & $-0.0059 \pm 0.0053$ \\
gaze-contingent @ 10 & $1.5377$ & $3.262$ & $0.9162$ & $-0.0365 \pm 0.0083$ \\
fixed-centre @ 40 & $1.5751$ & $3.306$ & $0.9182$ & $+0.0009 \pm 0.0038$ \\
fixed-centre @ 20 & $1.5627$ & $3.287$ & $0.9177$ & $-0.0115 \pm 0.0063$ \\
fixed-centre @ 10 & $1.5245$ & $3.228$ & $0.9160$ & $-0.0497 \pm 0.0087$ \\
\bottomrule
\end{tabular}
\end{table}


At the human cutoff the gaze-contingent arm measures $\Delta \IG = +0.0039 \pm 0.0032$~bits,
fold-paired. What this measures is a bound rather than a benefit. The two arms differ by at most
$0.0071$~bits per fixation, two standard errors above the fold-paired mean, under $1\%$ of the
$1.57$~bits the sharp control gains over the centerbias and a fifth of what foveation costs at
$f_{c0} = 10$.
The dose--response is monotone in strength, a lower cutoff meaning a stronger blur, from
$-0.0059 \pm 0.0053$ at $f_{c0} = 20$ to $-0.0365 \pm 0.0083$ at $10$.
The fixed-centre arms sit below the gaze-contingent ones at every cutoff (\tabref{phase17-arms}).
\figref{fold-paired} draws all six arm-versus-control differences with their ten fold values shown
individually.

\begin{figure}[tbp]
  \centering
  \includegraphics[width=\textwidth]{foveation_mit1003_initial/figs/fold_paired.png}
  \caption[The six arm-versus-control differences by fold]{The six differences of \tabref{phase17-arms} at fold level. Each circle is one fold, that fold's
    foveated read-out minus the control's over the images held out from both; the marker is the mean
    of the ten.}
  \label{fig:fold-paired}
\end{figure}

\subsection{The gaze-contingency contrast}
\label{subsec:results-gaze-contingency}

The gaze-contingent and the fixed-centre arm at each cutoff apply the identical blur profile, so the
difference between them is what following the gaze contributes with the amount of blur held constant
(\subsecref{arms}). That difference, the
gaze-contingent arm's $\Delta \IG$ minus the fixed-centre arm's at the same cutoff, is the
\emph{gaze-contingency term}.
It grows as the periphery degrades, from $+0.0030 \pm 0.0027$~bits at $f_{c0} = 40$ to
$+0.0057 \pm 0.0023$ at $20$ and $+0.0132 \pm 0.0040$ at $10$.

All three differences at a cutoff are taken against the same control on the same fixations, so they add,
the gaze-contingent arm's difference from the control being the fixed-centre arm's difference
plus the gaze-contingency term. Read that way, at $f_{c0} = 10$ the blur costs $0.0497$~bits and
following the gaze returns $0.0132$ of it, and at $20$ it costs $0.0115$ and returns $0.0057$.
Degrading the periphery costs the model, and putting the sharp region where the eye is gives part of
that cost back.

The gaze-contingent arm sometimes delivers the saccade target inside its sharp disc, which would explain
the term. That disc shrinks from $103.5$~px to nothing across the three cutoffs, and the
gaze-contingency term is largest at $f_{c0} = 10$, where the disc has closed.
Whatever produces the term requires the degradation to follow the gaze, and grows as the periphery is
degraded further.

\subsection{Where the cost falls}
\label{subsec:results-diagnostics}

The diagnostics of \subsecref{stats-plan} are computed on the validation split at the same reporting
epoch, pairing the two arms of each contrast on identical fixations.
The map-level diagnostics assign the strong-cutoff loss to spreading rather than to mislocation, the two
arms' AUC at $f_{c0} = 10$ differing by $0.0017$.
At $f_{c0} = 10$ the foveated arm's probability maps carry $0.043 \pm 0.012$~bits more entropy than the
control's, while its predicted mode does not move measurably, against a mode-to-fixation distance of
$191$~px in both arms.
At $f_{c0} = 20$ the same pattern is a quarter the size ($+0.011$~bits), and at $f_{c0} = 40$ the
difference is $-0.003$~bits.
A map spread evenly over the pixels of an MIT1003 stimulus would carry about $19.5$~bits and the
control's carry $17.0$, so $2.5$~bits is what concentrating the density achieves. Against that, none of these
differences reaches $2\%$.

The paired difference, which is a difference in log-likelihood as well
(\subsecref{metrics}), is binned by saccade amplitude with the bin edges at the sharp radii of
\eqnref{sharp-radius}, so every bin lies wholly inside or wholly outside an arm's sharp disc
(\subsecref{sharp-radius}; \figref{stratified-amplitude}).
At $f_{c0} = 40$ the profile shows no step at the $103.5$~px radius; pooling each fold's bins into the two
regions gives $+0.002 \pm 0.007$~bits inside the disc and $+0.005 \pm 0.005$ outside, a difference of
$+0.003 \pm 0.010$.
At $f_{c0} = 10$, where no disc exists, the cost deepens with amplitude and reaches
$-0.149 \pm 0.051$~bits beyond $500$~px.

Stratifying by fixation index tests the expectation that a gaze-contingent cost should build along the
scanpath (\figref{stratified-fixation-index}).
At $f_{c0} = 10$ it does the opposite, with the cost deepest over the first two scored fixations and
down to a fifth of that by the tail, while the two milder cutoffs show no comparable trend.
Neither the saccade amplitude nor the control's own level accounts for it. Mean saccade amplitude
\emph{rises} along the scanpath, from $149$ to $208$~px, and the cost at this cutoff deepens with
amplitude, so amplitude on its own would predict a cost that grows where the measured one shrinks. The
control's information gain falls from $1.80$ to about $1.45$~bits, so there is less to lose late in a
scanpath than early, and the cost shrinks faster than that. Taken as a share of the control's level it
runs from $2.9\%$ over the first scored fixation to $0.8\%$ over the ninth and later.
The rest is consistent with the earliest predictions having the
least history to condition on, so the image carries most of the prediction there, while later
predictions draw on a longer history, which the foveation leaves untouched.

\begin{figure}[tbp]
  \centering
  \includegraphics[width=\textwidth]{foveation_mit1003_initial/figs/stratified_amplitude.png}
  \caption[Information-gain difference by saccade amplitude]{The information-gain difference by
    saccade amplitude on the validation split, with each point and bar taken across the ten folds. The
    dashed lines are the sharp radii of \eqnref{sharp-radius}. The last bin is open-ended and holds
    every saccade beyond $500$~px.}
  \label{fig:stratified-amplitude}
\end{figure}

\begin{figure}[tbp]
  \centering
  \includegraphics[width=0.85\textwidth]{foveation_mit1003_initial/figs/stratified_fixation_index.png}
  \caption[Information-gain difference by fixation index]{The information-gain difference by fixation index on the validation split, counting from the
    first free fixation onwards (\subsecref{mit1003}); fixations nine and later are pooled.}
  \label{fig:stratified-fixation-index}
\end{figure}

\section{Per-image variation in model fit}
\label{sec:results-per-image}

This section reports how \dg{}'s per-image performance varies with the fixation entropy of
\subsecref{difficulty-proxies}.
The model is the pretrained \dg{}, and the sample is every MIT1003 stimulus under the protocol of
\subsecref{mit1003}.
\figref{per-image-scatter} plots the entropy against each of the four metrics, with Pearson's
product-moment $r$ and Spearman's rank $\rho$ over each panel.

\begin{figure}[tbp]
  \centering
  \includegraphics[width=0.88\textwidth]{per_image_diagnostic_initial_1003/scatter.png}
  \caption[Fixation entropy against model fit]{The fixation entropy of \eqnref{fix-entropy} against each of the four per-image metrics, one
    point per stimulus with a least-squares line, for the pretrained \dg{}; $r$ is Pearson's and
    $\rho$ Spearman's.}
  \label{fig:per-image-scatter}
\end{figure}

Fixation entropy runs against every metric, the direction \subsecref{difficulty-proxies} anticipates,
and most strongly against the log-likelihood.
Part of this is a property of the metric rather than of the model.
A likelihood rewards a concentrated density, so when the fifteen subjects looked at the same few places
any model that puts its mass there scores well, and when they looked everywhere none can.
The correlation therefore tracks how hard an image is to score on as much as how well \dg{} does on it.

Per-fixation information gain over the 1003 stimuli ranges from $-0.13$ to $3.96$~bits and AUC from $0.77$ to $0.98$.
\figref{extremes-high} shows the three highest stimuli by information gain and \figref{extremes-low}
the three lowest, with each stimulus's IG, NSS and AUC in the caption.
At the high end a region away from the image centre holds three-quarters of the subjects.
Stimulus $750$ is the one image where \dg{} scores below the centerbias.
The flower the subjects fixated sits at the image centre, where the centerbias already puts its mass,
so the prior that \dg{} is measured against is at its strongest there.

\begin{figure}[tbp]
  \centering
  \includegraphics[width=\linewidth]{consensus_panels/consensus_stim0128.png}\\[3pt]
  \includegraphics[width=\linewidth]{consensus_panels/consensus_stim0704.png}\\[3pt]
  \includegraphics[width=\linewidth]{consensus_panels/consensus_stim0528.png}
  \caption[The three highest stimuli by information gain]{The three stimuli with the highest per-fixation
    information gain, top to bottom: $128$ ($3.96$~bits, NSS $3.4$, AUC $0.931$), $704$ ($3.77$~bits,
    $5.7$, $0.968$) and $528$ ($3.64$~bits, $5.0$, $0.952$). In each row, left to right: every subject's
    scanpath on that stimulus; the consensus regions (Appendix~\ref{app:consensus}); \dg{}'s
    probability map given a single fixation at the image centre.}
  \label{fig:extremes-high}
\end{figure}

\begin{figure}[tbp]
  \centering
  \includegraphics[width=\linewidth]{consensus_panels/consensus_stim0750.png}\\[3pt]
  \includegraphics[width=\linewidth]{consensus_panels/consensus_stim0639.png}\\[3pt]
  \includegraphics[width=\linewidth]{consensus_panels/consensus_stim0489.png}
  \caption[The three lowest stimuli by information gain]{The three stimuli with the lowest per-fixation
    information gain, top to bottom: $750$ ($-0.13$~bits, NSS $2.6$, AUC $0.879$), $639$ ($0.23$~bits,
    $1.7$, $0.849$) and $489$ ($0.23$~bits, $1.8$, $0.846$). Panels as in \figref{extremes-high}.}
  \label{fig:extremes-low}
\end{figure}
